\subsubsection{Additionstheoreme}

Die Additionstheoreme geben Sinus und Kosinus einer Summe oder Differenz zweier
Winkel an.

\begin{TheoremBox}{Additionstheoreme}{fundamentals_trigonometric_addition_formulas}
  \begin{ClozeArea}[34]{8}
    \begin{align*}
      \sin(\alpha\pm\beta)
        &=\sin(\alpha)\cos(\beta)
          \pm\cos(\alpha)\sin(\beta),\\
      \cos(\alpha\pm\beta)
        &=\cos(\alpha)\cos(\beta)
          \mp\sin(\alpha)\sin(\beta).
    \end{align*}
  \end{ClozeArea}
\end{TheoremBox}

Besonders wichig ist der Spezialfall $\beta=\alpha$ für den man Doppelwinkelformeln erhält.

\begin{KeyPointBox}{Doppelwinkelformeln}{fundamentals_trigonometric_double_angle_formulas}
  \begin{ClozeArea}[28]{6}
    \begin{align*}
      \sin(2\alpha)&=2\sin(\alpha)\cos(\alpha),\\
      \cos(2\alpha)&=\cos^2(\alpha)-\sin^2(\alpha).
    \end{align*}
  \end{ClozeArea}
\end{KeyPointBox}

